Neural operators learn maps between function spaces, for example from the initial condition of a PDE to its solution at a later time, from a modest set of simulated pairs~\cite{kovachki2021neural}. Two families dominate the literature: the Fourier neural operator (FNO), which parametrises a kernel integral operator by a truncated Fourier multiplier~\cite{li2020fourier}, and DeepONet, which pairs a branch network acting on sensor values with a trunk network acting on query coordinates~\cite{lu2019deeponet}. A published comparison on 16 benchmarks reports that the two perform comparably in relatively simple settings, while the FNO deteriorates on complex geometries and noisy data~\cite{lu2021comprehensive}. PDEBench provides larger-scale data, including 1D Burgers, together with baselines~\cite{takamoto2022pdebench}. A frequently quoted property of the FNO is that, once trained, it can be applied on a finer grid than the one it was trained on (``zero-shot super-resolution'')~\cite{li2020fourier}; theory warns that this holds exactly only under conditions on aliasing and band-limitation~\cite{bartolucci2023representation}.

This paper asks how far these statements hold in the smallest setting one can reproduce on a laptop CPU in minutes. We generate 600 solution pairs of the viscous Burgers equation with a high-accuracy spectral solver, train four reimplemented models on 400 of them, and report the test relative $L_2$ error at the training resolution and at $2\times$ and $4\times$ finer grids. The MLP and CNN act as controls: neither is a neural operator, and the CNN, unlike the MLP, can at least be applied to any grid. We state four falsifiable hypotheses beforehand (Section~\ref{sec:setup}), and report each as supported, refuted or inconclusive.

\textbf{Contributions.} (i) A fully specified, five-seed protocol with the data generator, splits, tuning budget and runtimes. (ii) A measurement of zero-shot finer-grid error with explicit rules for how each fixed-grid model is applied, including a resampling variant. (iii) Sensitivity studies on Fourier modes, training-set size, training length and the grid channel. We claim no new method; all models are reimplementations and the study is small.
